\subsection{প্রতিরক্ষা}

\begin{learningbox}
এই অংশ শেষে শিক্ষার্থী পারবে:
\begin{itemize}
\item প্রতিরক্ষা ব্যবস্থায় ICT-এর ভূমিকা ব্যাখ্যা করতে।
\item radar, sensor, satellite surveillance, cyber security ও drone-এর ব্যবহার লিখতে।
\item জাতীয় নিরাপত্তায় তথ্য নিরাপত্তা ও যোগাযোগব্যবস্থার গুরুত্ব বুঝতে।
\end{itemize}
\end{learningbox}

\subsubsection{প্রতিরক্ষা ব্যবস্থায় আইসিটি}

প্রতিরক্ষা ব্যবস্থায় তথ্য সংগ্রহ, বিশ্লেষণ, যোগাযোগ, নজরদারি, command-control, simulation training এবং cyber security-তে ICT ব্যবহৃত হয়। আধুনিক প্রতিরক্ষা শুধু অস্ত্রের উপর নির্ভর করে না; দ্রুত, নির্ভুল ও নিরাপদ তথ্য ব্যবস্থাপনার উপরও নির্ভর করে।

\begin{definitionbox}{প্রতিরক্ষায় ICT}
প্রতিরক্ষা ব্যবস্থায় ICT হলো সেনা, নৌ, বিমান, সীমান্ত ও সাইবার নিরাপত্তায় computer, sensor, satellite, secure network, database, radar, drone ও software ব্যবহার করে তথ্য সংগ্রহ, যোগাযোগ ও সিদ্ধান্ত গ্রহণের প্রক্রিয়া।
\end{definitionbox}

\subsubsection{রাডার ও সেন্সর}

Radar radio wave ব্যবহার করে object-এর অবস্থান, দূরত্ব ও গতি নির্ণয় করে। বিমান, জাহাজ, missile বা weather system শনাক্ত করতে radar ব্যবহৃত হয়। Sensor আবার heat, sound, image, vibration, pressure বা movement detect করতে পারে।

\begin{keypointbox}{প্রতিরক্ষা sensor-এর কাজ}
\begin{itemize}
\item সীমান্তে movement detection।
\item aircraft বা ship tracking।
\item battlefield condition monitoring।
\item mine বা explosive detection।
\item surveillance camera ও thermal imaging।
\end{itemize}
\end{keypointbox}

\subsubsection{স্যাটেলাইট নজরদারি}

Satellite image ও communication system জাতীয় নিরাপত্তায় গুরুত্বপূর্ণ। সীমান্ত পর্যবেক্ষণ, দুর্যোগ-পরবর্তী ক্ষয়ক্ষতি দেখা, সমুদ্রসীমা পর্যবেক্ষণ, navigation এবং secure communication-এ satellite ব্যবহার করা যায়।

\subsubsection{সাইবার নিরাপত্তা}

আধুনিক রাষ্ট্রের বিদ্যুৎ, ব্যাংকিং, যোগাযোগ, সরকারি database, সামরিক network ও transport system computer network-এর উপর নির্ভরশীল। তাই cyber attack প্রতিরক্ষা ব্যবস্থার বড় ঝুঁকি। Firewall, encryption, authentication, intrusion detection, security monitoring ও digital forensics cyber security-তে ব্যবহৃত হয়।

\begin{mistakebox}{সাধারণ ভুল}
প্রতিরক্ষা মানেই শুধু যুদ্ধক্ষেত্র নয়। national database, communication network, power grid, banking system ও emergency service সুরক্ষাও আধুনিক প্রতিরক্ষার অংশ।
\end{mistakebox}

\subsubsection{ড্রোন ও রোবট}

Drone বা UAV দূর থেকে পরিচালিত বা autonomous flying device। নজরদারি, mapping, disaster response, border monitoring এবং কিছু সামরিক কাজে drone ব্যবহৃত হয়। Robot ব্যবহার করা যায় bomb disposal, hazardous material handling ও rescue operation-এ।

\begin{center}
\begin{tabular}{|p{0.28\textwidth}|p{0.58\textwidth}|}
\hline
\textbf{প্রযুক্তি} & \textbf{প্রতিরক্ষায় ব্যবহার} \\
\hline
Radar & aircraft, ship বা missile tracking \\
\hline
Satellite & communication, navigation, surveillance \\
\hline
Cyber security & network ও data protection \\
\hline
Drone & remote surveillance ও mapping \\
\hline
Robot & bomb disposal ও risky task \\
\hline
\end{tabular}
\end{center}

\begin{boardbox}{বোর্ড ফোকাস}
প্রতিরক্ষায় ICT-এর প্রশ্নে radar, satellite, cyber security, drone এবং secure communication—এই পাঁচটি keyword ব্যবহার করো।
\end{boardbox}

\begin{practicebox}
\textbf{MCQ}
\begin{enumerate}
\item Radar কোন wave ব্যবহার করে object detect করে?\\
ক. radio wave \quad খ. water wave \quad গ. sound only \quad ঘ. paper signal
\item Network attack থেকে system রক্ষার ক্ষেত্র কোনটি?\\
ক. Cyber security \quad খ. Word processing \quad গ. Desktop publishing \quad ঘ. Formatting
\end{enumerate}

\textbf{সংক্ষিপ্ত প্রশ্ন}
\begin{enumerate}
\item প্রতিরক্ষায় satellite-এর দুটি ব্যবহার লেখ।
\item Cyber security কেন জাতীয় নিরাপত্তার অংশ?
\end{enumerate}
\end{practicebox}

